\section{Introduction}
\label{sec:intro}

A child hears no more than about 100 million words by early adolescence. A
language model consumes several orders of magnitude more before it is
called trained. The BabyLM Challenge asks what this gap conceals: how much
linguistic competence can pretraining extract from a corpus a child could
plausibly have heard? \citep{warstadt2023findings,charpentier2025findings}
We compete in the Strict-Small track of the 2026 edition
\citep{choshen2026babylm}, the tightest form of the question: at most 10
million words of text, with the corpus traversed at most ten times.

Most entries respond by adding signal. New objectives, curricula,
architectures, tokenizers, and distilled teachers arrive every year
\citep{hu2024findings,charpentier2025findings}. We start from a different
question: before adding anything, when does the existing budget do its work?
Training a \rParamsM{}-parameter GPT-BERT hybrid
\citep{charpentier2024gptbert} under the track rules, we find that
downstream competence saturates by roughly \rPeakCkpt{} words seen, seven of
the ten permitted epochs, while training loss keeps falling. Those last
epochs purchase a lower loss, not new competence. A per-phenomenon breakdown
adds a second fact: a small cluster of long-distance dependencies is
acquired last and stays fragile to the end.

The diagnosis dictates what a useful intervention must look like. It should
arrive after saturation, where its gradient no longer competes with language
modeling for capacity. It must be anchored, because fragile structure is
what strong new gradients destroy first. And it should share the functional
form of the evaluation, which for the core benchmarks is a likelihood
comparison between minimally different strings. We call these the timing,
preservation, and alignment principles.

Their instantiation is \MPOlong{} (\MPO{}): a short preference-optimization
phase \citep{rafailov2023dpo} that restarts from the saturation checkpoint
and teaches the model to prefer real corpus sentences over length-preserving
corruptions of them, with a frozen reference penalizing drift
(Figure~\ref{fig:abstract}). The principles are tested rather than assumed.
A controlled study of seven in-training interventions drawn from BabyLM
entries and the broader pretraining literature confirms both predictions:
none improves on the baseline, and the unanchored ones damage the fragile
cluster first. \MPO{} improves five of the seven official task columns and concedes
under a point on the rest, is alone in repairing the fragile cluster,
and matches the fully trained baseline after \rDpoEpochs{} rather than
\rEpochs{} epochs of exposure; our \MPO{} entry places \rRankDpo{} on the
track leaderboard, the fully trained baseline \rRankBase{}.

Beyond the method, the paper contributes the saturation diagnosis, cheap
to run and decisive about where post-training signal belongs, and the
controlled study of why published gains fail to transfer here.
